
\documentclass[a4paper,10pt]{report}

%\usepackage[landscape]{geometry}
\usepackage[cm-default]{fontspec}
\setromanfont{FreeSerif}
\setsansfont{FreeSans}
\setmonofont{FreeMono}
\usepackage[utf8x]{inputenc}
\usepackage{fontspec}
\usepackage{xunicode}
\usepackage{xltxtra}
\usepackage{xgreek}
\usepackage{amsmath}
\usepackage{unicode-math}
\usepackage{ulem}
\usepackage{color}
\usepackage{verbatim}
\usepackage{nopageno}
\usepackage{graphicx}
\usepackage{textpos}
\setlength{\TPHorizModule}{1cm}
\setlength{\TPVertModule}{1cm}
%\usepackage[colorgrid,texcoord]{eso-pic}
\usepackage[outline]{contour}
\usepackage{wrapfig}
\usepackage{url}
\usepackage{color}
\usepackage{tikz}
\usepackage{fancybox,fancyhdr}
\usepackage{subfigure}
\usepackage{pstricks}
\usepackage{epsfig}
\usepackage{multicol}
\usepackage{listings}
\usepackage{enumerate}
\usepackage{hyperref}
\hypersetup{
  bookmarks=true,
  bookmarksopen=true,
  pdfborder=false,
  pdfpagemode=UseNone,
  raiselinks=true,
  pdfhighlight={/P},
  colorlinks,
  citecolor=black,
  filecolor=black,
  linkcolor=black,
  urlcolor=black
}


\usepackage{fontspec}
\usepackage{xunicode}
\usepackage{xltxtra}


% Margins
%\setlength{\textwidth}{24cm}

%\setlength{\voffset}{0in}
%\setlength{\textheight}{6.8in}
% Colors
\definecolor{cbrown}{rgb}{0.49,0.24,0.07}
\definecolor{cmpez}{rgb}{0.92,0.88,0.79}
\definecolor{cmpez2}{rgb}{0.62,0.33,0.34}
\definecolor{cwhite}{rgb}{1,1,1}
\definecolor{cblack}{rgb}{0,0,0}
\definecolor{cred}{rgb}{0.9,0.15,0.15}
\definecolor{clblue}{rgb}{0.57,0.85,0.97}
\definecolor{corange}{rgb}{0.86,0.49,0.18}
\definecolor{clorange}{rgb}{0.95,0.84,0.61}
\definecolor{cyellow}{rgb}{1,0.95,0.16}
\definecolor{cgreen}{rgb}{0.69,0.8,0.31}
\definecolor{clbrown}{rgb}{0.78,0.67,0.43}
\definecolor{cmagenta}{rgb}{0.79,0.05,0.54}
\definecolor{cgray}{rgb}{0.76,0.73,0.63}
\definecolor{bluesite}{rgb}{0.05,0.43,0.67}
\definecolor{redsite}{rgb}{0.90,0.15,0.15}
\definecolor{backsite}{rgb}{0.07,0.07,0.07}
% ----------- from Costas -----------------------------

%---------------   Main Fonts    -----------------%   

% \setmainfont[Mapping=tex-text]{Calibri}
% \setmainfont[Mapping=tex-text]{Times New Roman}
%\setmainfont[Mapping=tex-text]{Droid Serif}
%\setmainfont[Mapping=tex-text]{Cambria}
%\setmainfont[Mapping=tex-text]{cm-unicode}
% \setmainfont[Mapping=tex-text]{Gentium}
% \setmainfont[Mapping=tex-text]{GFS Didot}
% \setmainfont[Mapping=tex-text]{Comic Sans MS}
 \setmainfont[Mapping=tex-text]{Ubuntu}
% \setmainfont[Mapping=tex-text]{Myriad Pro}
% \setmainfont[Mapping=tex-text]{CMU Concrete}
%\setmainfont[Mapping=tex-text]{DejaVu Sans}
%\setmainfont[Mapping=tex-text]{KerkisSans}
%\setmainfont[Mapping=tex-text]{KerkisCaligraphic}
% \setmainfont[Mapping=tex-text]{Segoe Print}
% \setmainfont[Mapping=tex-text]{Gabriola}

%------------------------------------------------%

% ------------ Mathematics Fonts ----------------%
%\setmathfont{Asana-Math.ttf}
% -----------------------------------------------%
% Misc
\author{Κασωτάκη Ε. - Σμαραγδάκης Κ.}
\title{www.math24.gr}
% ---------------- Formation --------------------%


\setlength\topmargin{-0.7cm}
\addtolength\topmargin{-\headheight}
\addtolength\topmargin{-\headsep}
\setlength\textheight{26cm}
\setlength\oddsidemargin{-0.54cm}
\setlength\evensidemargin{-0.54cm}
%\setlength\marginparwidth{1.5in}
\setlength\textwidth{17cm}

\RequirePackage[avantgarde]{quotchap}
\renewcommand\chapterheadstartvskip{\vspace*{0\baselineskip}}
\RequirePackage[calcwidth]{titlesec}
\titleformat{\section}[hang]{\bfseries}
{\Large\thesection}{12pt}{\Large}[{\titlerule[0.9pt]}]
%--------------------------------------------------%

%\usepackage{draftwatermark}
%\SetWatermarkText{www.math24.gr}
%\SetWatermarkLightness{0.9}
%\SetWatermarkScale{0.6}

%--------------------------------------------------%


\usepackage{xcolor}
\usepackage{amsthm}
\usepackage{framed}
\usepackage{parskip}

\colorlet{shadecolor}{bluesite!20}

\newtheorem*{orismos}{Ορισμός}

\newenvironment{notation}
  {\begin{shaded}\begin{theorem}}
  {\end{theorem}\end{shaded}}




% Document begins
\begin{document}
\pagestyle{fancy}
\fancyhead{}
\fancyfoot{}
\renewcommand{\headrulewidth}{0pt}
\renewcommand{\footrulewidth}{0pt}


\fancyhead[LO,LE]{
 %\textblockcolor{backsite}
 %\begin{textblock}{5}(-2,-0.55)
  %\rule{0cm}{1cm}
 %\end{textblock}
 \textblockcolor{bluesite}
 \begin{textblock}{5}(-1.5,-0.55)
  \rule{0cm}{1cm}
 \end{textblock}
 %\textblockcolor{bluesite}
 %\begin{textblock}{14}(5.5,-0.55)
 % \rule{0cm}{1cm}
 %\end{textblock}
 \begin{textblock}{0}(-1,-0.25)
 \color{cwhite} \begin{Large}www.math24.gr\end{Large}
 \end{textblock}
\begin{textblock}{0}(16,-0.4)
 \color{cwhite} \includegraphics[height=1.5cm]{math24_logo.png}
 \end{textblock}
\textblockcolor{white}
\begin{textblock}{17}(-1,28)
 \color{backsite} \begin{small}Copyright \textcopyright 2018, Κασωτάκη Ε.(ikasotaki@gmail.com) - Σμαραγδάκης Κ.(kesmarag@gmail.com)\end{small}
 \end{textblock}
}
 
\begin{shaded}
\begin{center}
\huge \textbf{Φυλλάδιο Ασκήσεων}\\
%Πρόσθεση ρητών αριθμών
\end{center} 
\textbf{Μαθηματικά Γ' Γυμνασίου} \hfill \textbf{Ημερομηνία Παράδοσης : \hspace{2em} }
\subsection*{Ονοματεπώνυμο :\hfill  \hspace{5em}}
\end{shaded}
\vspace{2em}
\begin{itemize}
 \item πρόσθεση μονωνύμων
 \item πολλαπλασιασμός μονωνύμων
 \item διαίρεση μονωνύμων
\end{itemize}

\section*{Θεωρία - Πρόσθεση μονωνύμων \hfill \small{}}
\begin{itemize}
 \item Αν \textbf{προσθέτουμε} όμοια μονώνυμα τότε το άθροισμα τους είναι μονώνυμο όμοιο με αυτά και έχει συντελεστή 
       το άθροισμα των συντελεστών τους (όπως προκύπτει με χρήση της επιμεριστικής ιδιότητας).\\
       \textbf{Π.χ} $3xy^{2}+4xy^{2}=(3+4)xy^{2}=7xy^{2}$
 \item Αν \textbf{προσθέτουμε} όχι όμοια μονώνυμα τότε το άθροισμά τους δεν είναι μονώνυμο.\\
       \textbf{Π.χ} $(7x)+(6x^{2})=7x+6x^{2}$ το οποίο δεν είναι μονώνυμο
\end{itemize}


\section*{Άσκηση 1  \hfill \small{ }}
Να κάνετε τις παρακάτω πράξεις:
\begin{enumerate}[i)]
 \item $3x^{2}+5x^{2}$
 \item $-2y^{3}x+4y^{3}x$
 \item $6x^{2}y^{3}-2x^{2}y^{3}+10y^{3}x^{2}$
 \item $2x^{5}+6x^{5}+xy$
 \item $7x^{4}y+8x^{3}y+3x^{4}y-3x^{3}y$
\end{enumerate}

\section*{Θεωρία - Πολλαπλασιασμός και Διαίρεση μονωνύμων \hfill \small{}}
\begin{itemize}
 \item Όταν \textbf{πολλαπλασιάζουμε} μονώνυμα τότε το γινόμενό τους είναι μονώνυμο με συντελεστή το 
       γινόμενο των συντελεστών τους και κύριο μέρος το γινόμενο όλων των μεταβλητών τους με εκθέτη 
       κάθε μεταβλητής το άθροισμα των εκθετών της.\\
       \textbf{Π.χ} $(3x^{3}y^{2})\cdot (5xy^{4}) =15x^{4}y^{6} $
 \item Για να \textbf{διαιρέσουμε} δύο μονώνυμα, πολλαπλασιάζουμε το διαιρετέο με τον αντίστροφο του 
       διαιρέτη (αλλά δεν προκύπτει πάντα μονώνυμο).\\
       \textbf{Π.χ} $ (8x^{4}y^{3}):(4x^{2}y^{2}) = 8x^{4}y^{3}\cdot \dfrac{1}{4x^{2}y^{2}}=
        \dfrac{8x^{4}y^{3}}{4x^{2}y^{2}}=2x^{2}y$
\end{itemize}



\section*{Άσκηση 2  \hfill \small{ }}
Να κάνετε τις παρακάτω πράξεις:
\begin{enumerate}[i)]
 \item $2xy^{6} \cdot 5x^{3}y$
 \item $-4x^{2}y \cdot 5y^{3}$
 \item $(-\dfrac{1}{2}x^{3}y^{3})\cdot (-8x^{4}y^{4})$
 \item $5xy \cdot 5x^{2}y^{2}$ 
 \item $\sqrt{3}x^{3}y^{3} \cdot \sqrt{3}x^{2}y^{4}$
\end{enumerate}









\section*{Άσκηση 3  \hfill \small{ }}
Να κάνετε τις παρακάτω πράξεις:
\begin{enumerate}[i)]
 \item $10x^{3}y : 5x^{2}y$
 \item $15x^{4}y : 20x^{2}y$
 \item $14x^{2}y^{5} : 7x^{4}y^{3}$
 \item $(-12y^{3}) : (-4x^{2}y^{2})$
 \item $6x^{4}y^{4} : (-5x^{5}y^{5})$
\end{enumerate}



\section*{Άσκηση 4  \hfill \small{ }}
Να κάνετε τις παρακάτω πράξεις:
\begin{enumerate}[i)]
 \item $\dfrac{1}{3}xy^{2}+\dfrac{2}{5}xy^{2}$
 \item $\dfrac{1}{3}xy^{2}-\dfrac{2}{5}xy^{2}$
 \item $(\dfrac{1}{3}xy^{2}) \cdot (\dfrac{2}{5}xy^{2})$
 \item $(\dfrac{1}{3}xy^{2}) : (\dfrac{2}{5}xy^{2})$
\end{enumerate}





 









%\let\thefootnote\relax\footnotetext{Επιμέλεια : Κασωτάκη Ειρήνη (ikasotaki@gmail.com),    Σμαραγδάκης Κώστας (kesmarag@gmail.com)} 


\end{document}
